% GENERATED FILE. Edit canonical lessons, then run npm run generate:books.
% canonical-source-hash: fnv1a64:ef12832679f43844
% canonical-lessons: GE-C405-kofferraum, GE-C405-motorhaube, GE-C405-kratzer, GE-C405-delle, GE-C405-verbeulen, GE-R405-review-398-405

\chapter{Kratzer und Dellen}
\label{ch:ge-kratzer-und-dellen}
\hlchaptermodality{405}

\begin{chapteropening}
\textbf{By the end of this chapter:} \emph{I can describe damage to a car: the boot, the bonnet, a scratch, a dent and denting something.}

Complete the last lesson of chapter 405: I can describe damage to a car: the boot, the bonnet, a scratch, a dent and denting something.
\end{chapteropening}

\section[der Kofferraum]{der Kofferraum — the boot (of a car; US: trunk)}

\label{lesson:GE-C405-kofferraum}



\begin{quote}
Before the new one: say the German for the jack, then the German for the warning triangle.
\end{quote}

\subsection*{You'll want to know: der Kofferraum}
\textbf{der Kofferraum} — \textquotedblleft{}the boot (of a car; US: trunk)\textquotedblright{}.

\begin{quote}
Die Koffer passen nicht alle in den Kofferraum.
\end{quote}

\begin{quote}
\emph{The suitcases do not all fit in the boot.}
\end{quote}

Plural: \emph{die Kofferräume}.

\begin{cousinweb}
\emph{Kofferraum} is built from \emph{Koffer} 'suitcase', borrowed from French \emph{coffre} 'chest' (EtymDB), and \emph{Raum} 'space', inherited from Old High German \emph{rūm} (EtymDB).
\end{cousinweb}

\begin{grammarlens}[title={What you've built}]
One more word for telling, explaining and describing at B1.
\end{grammarlens}

\subsection*{Your turn}
\begin{itemize}
\raggedright
  \item \emph{Say it:} \emph{der Kofferraum}
  \item \emph{Say it:} \emph{der Kofferraum}, once more
  \item \emph{Say it:} a sentence of your own with \emph{der Kofferraum}
  \item \emph{Recall:} say the German for the jack, then the German for the warning triangle, then say \emph{der Kofferraum} again
\end{itemize}

\begin{tcolorbox}[breakable,colback=teal!4,colframe=teal!35!black,title={Before you move on}]
What does der Kofferraum mean? (\textquotedblleft{}The boot (of a car; US: trunk)\textquotedblright{}.) Say it once more.
\end{tcolorbox}
\section[die Motorhaube]{die Motorhaube — the bonnet (of a car; US: hood)}

\label{lesson:GE-C405-motorhaube}



\begin{quote}
Before the new one: say the German for the warning triangle, then the German for the boot.
\end{quote}

\subsection*{You'll want to know: die Motorhaube}
\textbf{die Motorhaube} — \textquotedblleft{}the bonnet (of a car; US: hood)\textquotedblright{}.

\begin{quote}
Der Mechaniker öffnet die Motorhaube und schaut lange hinein.
\end{quote}

\begin{quote}
\emph{The mechanic opens the bonnet and looks inside for a long time.}
\end{quote}

Plural: \emph{die Motorhauben}.

\begin{cousinweb}
\emph{Motorhaube} is built from \emph{Motor} 'engine', from Latin \emph{mōtor} 'mover' (EtymDB; Lewis \& Short), and \emph{Haube} 'cap, hood', an inherited Germanic word for a head covering.
\end{cousinweb}

\begin{grammarlens}[title={What you've built}]
One more word for telling, explaining and describing at B1.
\end{grammarlens}

\subsection*{Your turn}
\begin{itemize}
\raggedright
  \item \emph{Say it:} \emph{die Motorhaube}
  \item \emph{Say it:} \emph{die Motorhaube}, once more
  \item \emph{Say it:} a sentence of your own with \emph{die Motorhaube}
  \item \emph{Recall:} say the German for the warning triangle, then the German for the boot, then say \emph{die Motorhaube} again
\end{itemize}

\begin{tcolorbox}[breakable,colback=teal!4,colframe=teal!35!black,title={Before you move on}]
What does die Motorhaube mean? (\textquotedblleft{}The bonnet (of a car; US: hood)\textquotedblright{}.) Say it once more.
\end{tcolorbox}
\section[der Kratzer]{der Kratzer — the scratch}

\label{lesson:GE-C405-kratzer}



\begin{quote}
Before the new one: say the German for the boot, then the German for the bonnet.
\end{quote}

\subsection*{You'll want to know: der Kratzer}
\textbf{der Kratzer} — \textquotedblleft{}the scratch\textquotedblright{}.

\begin{quote}
Auf der Tür ist ein langer Kratzer.
\end{quote}

\begin{quote}
\emph{There is a long scratch on the door.}
\end{quote}

Same form in the plural.

\begin{cousinweb}
\emph{Kratzer} is built on the verb \emph{kratzen} 'to scratch' (EtymDB), which is inherited from Old High German \emph{krazzōn}.
\end{cousinweb}

\begin{grammarlens}[title={What you've built}]
One more word for telling, explaining and describing at B1.
\end{grammarlens}

\subsection*{Your turn}
\begin{itemize}
\raggedright
  \item \emph{Say it:} \emph{der Kratzer}
  \item \emph{Say it:} \emph{der Kratzer}, once more
  \item \emph{Say it:} a sentence of your own with \emph{der Kratzer}
  \item \emph{Recall:} say the German for the boot, then the German for the bonnet, then say \emph{der Kratzer} again
\end{itemize}

\begin{tcolorbox}[breakable,colback=teal!4,colframe=teal!35!black,title={Before you move on}]
What does der Kratzer mean? (\textquotedblleft{}The scratch\textquotedblright{}.) Say it once more.
\end{tcolorbox}
\section[die Delle]{die Delle — the dent}

\label{lesson:GE-C405-delle}



\begin{quote}
Before the new one: say the German for the bonnet, then the German for the scratch.
\end{quote}

\subsection*{You'll want to know: die Delle}
\textbf{die Delle} — \textquotedblleft{}the dent\textquotedblright{}.

\begin{quote}
Nach dem Parken hat das Auto eine kleine Delle.
\end{quote}

\begin{quote}
\emph{After parking, the car has a small dent.}
\end{quote}

Plural: \emph{die Dellen}. Also \emph{die Beule}.

\begin{cousinweb}
\emph{Delle} is inherited from Middle High German \emph{telle} 'a hollow, a dip in the ground' (EtymDB).
\end{cousinweb}

\begin{grammarlens}[title={What you've built}]
One more word for telling, explaining and describing at B1.
\end{grammarlens}

\subsection*{Your turn}
\begin{itemize}
\raggedright
  \item \emph{Say it:} \emph{die Delle}
  \item \emph{Say it:} \emph{die Delle}, once more
  \item \emph{Say it:} a sentence of your own with \emph{die Delle}
  \item \emph{Recall:} say the German for the bonnet, then the German for the scratch, then say \emph{die Delle} again
\end{itemize}

\begin{tcolorbox}[breakable,colback=teal!4,colframe=teal!35!black,title={Before you move on}]
What does die Delle mean? (\textquotedblleft{}The dent\textquotedblright{}.) Say it once more.
\end{tcolorbox}
\section[verbeulen]{verbeulen — to dent}

\label{lesson:GE-C405-verbeulen}



\begin{quote}
Before the new one: say the German for the scratch, then the German for the dent.
\end{quote}

\subsection*{You'll want to know: verbeulen}
\textbf{verbeulen} — \textquotedblleft{}to dent\textquotedblright{}.

\begin{quote}
Ein Ball hat die Tür von unserem Auto verbeult.
\end{quote}

\begin{quote}
\emph{A ball dented the door of our car.}
\end{quote}

Often as an adjective: \emph{ein verbeultes Auto} — a dented car.

\begin{cousinweb}
\emph{verbeulen} is built from \emph{ver-} and \emph{Beule} 'bump, dent', which is inherited through Middle High German \emph{biule} from Old High German \emph{būlla} (EtymDB).
\end{cousinweb}

\begin{grammarlens}[title={What you've built}]
One more word for telling, explaining and describing at B1.
\end{grammarlens}

\subsection*{Your turn}
\begin{itemize}
\raggedright
  \item \emph{Say it:} \emph{verbeulen}
  \item \emph{Say it:} \emph{verbeulen}, once more
  \item \emph{Say it:} a sentence of your own with \emph{verbeulen}
  \item \emph{Recall:} say the German for the scratch, then the German for the dent, then say \emph{verbeulen} again
\end{itemize}

\begin{tcolorbox}[breakable,colback=teal!4,colframe=teal!35!black,title={Before you move on}]
What does verbeulen mean? (\textquotedblleft{}To dent\textquotedblright{}.) Say it once more.
\end{tcolorbox}
\section[(dialogue)]{Auto und Straße}

\label{lesson:GE-R405-review-398-405}



\begin{quote}
Say the German for the country road and the exit.
\end{quote}

\subsection*{Your turn}
Say these aloud:

\begin{itemize}
\raggedright
  \item the country road, the exit, the slip road, the roundabout, to turn off
  \item the steering wheel, to steer, to brake, the brake, to change gear
  \item the headlight, the windscreen, the windscreen wiper, the rear-view mirror, to indicate
  \item the horn, to hoot, the seat belt, to buckle up, the front-seat passenger
  \item to fill up, to fill the tank right up, petrol, the petrol pump, the oil change
  \item to start, the battery, the jump start, the tow truck, the mechanic
  \item the tyre, the flat tyre, to pump up, the jack, the warning triangle
  \item the boot, the bonnet, the scratch, the dent, to dent
\end{itemize}

\begin{tcolorbox}[breakable,colback=teal!4,colframe=teal!35!black,title={Before you move on}]
The country road? (\textbf{die Landstraße}.) The exit? (\textbf{die Ausfahrt}.) The slip road? (\textbf{die Auffahrt}.) The roundabout? (\textbf{der Kreisverkehr}.) To turn off? (\textbf{abbiegen}.) The steering wheel? (\textbf{das Lenkrad}.) To steer? (\textbf{lenken}.) To brake? (\textbf{bremsen}.) The brake? (\textbf{die Bremse}.) To change gear? (\textbf{schalten}.) The headlight? (\textbf{der Scheinwerfer}.) The windscreen? (\textbf{die Windschutzscheibe}.) The windscreen wiper? (\textbf{der Scheibenwischer}.) The rear-view mirror? (\textbf{der Rückspiegel}.) To indicate? (\textbf{blinken}.) The horn? (\textbf{die Hupe}.) To hoot? (\textbf{hupen}.) The seat belt? (\textbf{der Sicherheitsgurt}.) To buckle up? (\textbf{anschnallen}.) The front-seat passenger? (\textbf{der Beifahrer}.) To fill up? (\textbf{tanken}.) To fill the tank right up? (\textbf{volltanken}.) Petrol? (\textbf{das Benzin}.) The petrol pump? (\textbf{die Zapfsäule}.) The oil change? (\textbf{der Ölwechsel}.) To start? (\textbf{anspringen}.) The battery? (\textbf{die Batterie}.) The jump start? (\textbf{die Starthilfe}.) The tow truck? (\textbf{der Abschleppwagen}.) The mechanic? (\textbf{der Mechaniker}.) The tyre? (\textbf{der Reifen}.) The flat tyre? (\textbf{der Platten}.) To pump up? (\textbf{aufpumpen}.) The jack? (\textbf{der Wagenheber}.) The warning triangle? (\textbf{das Warndreieck}.) The boot? (\textbf{der Kofferraum}.) The bonnet? (\textbf{die Motorhaube}.) The scratch? (\textbf{der Kratzer}.) The dent? (\textbf{die Delle}.) To dent? (\textbf{verbeulen}.)
\end{tcolorbox}
